% ЗГЕНЕРОВАНО scripts/build_format_v2_chapter.py — не редагувати вручну.
% Джерело: books/ashtavakra-gita/poetychnyi/rozdil-1.md

\poemverse{1.1}{Як знання до людини приходять? \\ Як свобода в ній настає? \\ Як безсторонність вона пізнає? \\ Скажи мені, пане, — усе.}
\poemverse{1.2}{Якщо волі, сину, ти хочеш — \\ відкинь чуттєве, мов отруту. \\ Терпимість, щирість, співчуття, \\ вдоволеність, правдивість \\ шануй, мов нектар, повсякчас.}
\poemverse{1.3}{Ти — не земля, не вода, не вогонь, \\ не вітер, не небо, і не з них. \\ Заради волі пізнай Атман — \\ свідка стихій, чия природа — cit.}
\poemverse{1.4}{Якщо, від тіла відірвавшись, \\ у cit спочинеш і стоїш, — \\ уже тепер ти будеш щасливий, \\ умиротворений, без пут.}
\poemverse{1.5}{Ти не брахман і не інший стан, \\ не учень і не відречений, \\ і око тебе не бачить. \\ Ти без прив'язності, без форми, \\ свідок усього. Будь щасливий.}
\poemverse{1.6}{Дхарма й адхарма, щастя й мука — \\ усе це в умі, не твоє, всепроникний. \\ Ти не дієць і не споживач — \\ ти вільний завжди.}
\poemverse{1.7}{Ти єдиний бачитель усього \\ і завжди майже вільний. \\ Ось у чому твоя уза: \\ ти бачиш іншого бачителя, крім себе.}
\poemverse{1.8}{«Я діяч» — це самомніння, \\ укус великого чорного змія. \\ Випий нектар упевненості: \\ «я не діяч» — і будь щасливий.}
\poemverse{1.9}{«Я — єдине чисте прозріння»: \\ цим вогнем упевненості \\ спали хащу незнання — \\ будь без жалю, будь щасливий.}
\poemverse{1.10}{Там, де все це сяє — \\ вигадане, як змія в мотузці, — \\ ти є те прозріння, \\ що є радість і найвища радість. \\ Ходи щасливо.}
\poemverse{1.11}{Хто вважає себе вільним — вільний; \\ хто вважає себе зв'язаним — зв'язаний. \\ Що тут іще сказати? \\ Яке переконання — така й доля.}
\poemverse{1.12}{Без прив'язності, без жадання, втихлий — \\ через оману він видається тим, хто в сансарі.}
\poemverse{1.13}{Осягай Атман — незмінне, \\ недвоїсте усвідомлення. \\ Відкинувши оману «я — лише відблиск», \\ відкинь і стан зовнішній, і внутрішній.}
\poemverse{1.14}{Ти довго зв'язаний, синку, \\ мотузкою приписування тіла. \\ Скажи: «я — усвідомлення» — \\ і відсічи те мечем знання. \\ Будь щасливий.}
\poemverse{1.15}{Ти без прив'язності, без діяння, \\ самосяйний, незаплямлений. \\ Ось саме це твої пута: \\ те, що ти вправляєшся в самадгі.}
\poemverse{1.16}{Цей всесвіт просякнутий тобою \\ і справді нанизаний на тебе. \\ Ти маєш природу чистого й просвітленого — \\ не йди до дріб'язковості ума.}
\poemverse{1.17}{Без очікування, без видозмін, \\ без тягаря, з прохолодним нутром, \\ із незміренною думкою, незбурений — \\ будь тим, чия схильність є сама лише свідомість.}
\poemverse{1.18}{Знай: те, що має образ, — несправжнє; \\ а те, що без образу — непохитне. \\ Цим настановленням про суть \\ немає народження знову.}
\poemverse{1.19}{Як образ у дзеркалі — \\ він усередині й навколо нього; \\ так само і в цьому тілі — \\ всередині й навколо верховний володар.}
\poemverse{1.20}{Як один усюдисущий простір — \\ зовні й усередині глечика, \\ так повсякчас і без перерви \\ брахман у всій сукупності істот.}
